% Options for packages loaded elsewhere
\PassOptionsToPackage{unicode}{hyperref}
\PassOptionsToPackage{hyphens}{url}
%
\documentclass[
  12pt,
]{article}
\usepackage{amsmath,amssymb}
\usepackage{iftex}
\ifPDFTeX
  \usepackage[T1]{fontenc}
  \usepackage[utf8]{inputenc}
  \usepackage{textcomp} % provide euro and other symbols
\else % if luatex or xetex
  \usepackage{unicode-math} % this also loads fontspec
  \defaultfontfeatures{Scale=MatchLowercase}
  \defaultfontfeatures[\rmfamily]{Ligatures=TeX,Scale=1}
\fi
\usepackage{lmodern}
\ifPDFTeX\else
  % xetex/luatex font selection
\fi
% Use upquote if available, for straight quotes in verbatim environments
\IfFileExists{upquote.sty}{\usepackage{upquote}}{}
\IfFileExists{microtype.sty}{% use microtype if available
  \usepackage[]{microtype}
  \UseMicrotypeSet[protrusion]{basicmath} % disable protrusion for tt fonts
}{}
\makeatletter
\@ifundefined{KOMAClassName}{% if non-KOMA class
  \IfFileExists{parskip.sty}{%
    \usepackage{parskip}
  }{% else
    \setlength{\parindent}{0pt}
    \setlength{\parskip}{6pt plus 2pt minus 1pt}}
}{% if KOMA class
  \KOMAoptions{parskip=half}}
\makeatother
\usepackage{xcolor}
\usepackage[margin=1in]{geometry}
\usepackage{longtable,booktabs,array}
\usepackage{calc} % for calculating minipage widths
% Correct order of tables after \paragraph or \subparagraph
\usepackage{etoolbox}
\makeatletter
\patchcmd\longtable{\par}{\if@noskipsec\mbox{}\fi\par}{}{}
\makeatother
% Allow footnotes in longtable head/foot
\IfFileExists{footnotehyper.sty}{\usepackage{footnotehyper}}{\usepackage{footnote}}
\makesavenoteenv{longtable}
\usepackage{graphicx}
\makeatletter
\def\maxwidth{\ifdim\Gin@nat@width>\linewidth\linewidth\else\Gin@nat@width\fi}
\def\maxheight{\ifdim\Gin@nat@height>\textheight\textheight\else\Gin@nat@height\fi}
\makeatother
% Scale images if necessary, so that they will not overflow the page
% margins by default, and it is still possible to overwrite the defaults
% using explicit options in \includegraphics[width, height, ...]{}
\setkeys{Gin}{width=\maxwidth,height=\maxheight,keepaspectratio}
% Set default figure placement to htbp
\makeatletter
\def\fps@figure{htbp}
\makeatother
\setlength{\emergencystretch}{3em} % prevent overfull lines
\providecommand{\tightlist}{%
  \setlength{\itemsep}{0pt}\setlength{\parskip}{0pt}}
\setcounter{secnumdepth}{5}
% definitions for citeproc citations
\NewDocumentCommand\citeproctext{}{}
\NewDocumentCommand\citeproc{mm}{%
  \begingroup\def\citeproctext{#2}\cite{#1}\endgroup}
\makeatletter
 % allow citations to break across lines
 \let\@cite@ofmt\@firstofone
 % avoid brackets around text for \cite:
 \def\@biblabel#1{}
 \def\@cite#1#2{{#1\if@tempswa , #2\fi}}
\makeatother
\newlength{\cslhangindent}
\setlength{\cslhangindent}{1.5em}
\newlength{\csllabelwidth}
\setlength{\csllabelwidth}{3em}
\newenvironment{CSLReferences}[2] % #1 hanging-indent, #2 entry-spacing
 {\begin{list}{}{%
  \setlength{\itemindent}{0pt}
  \setlength{\leftmargin}{0pt}
  \setlength{\parsep}{0pt}
  % turn on hanging indent if param 1 is 1
  \ifodd #1
   \setlength{\leftmargin}{\cslhangindent}
   \setlength{\itemindent}{-1\cslhangindent}
  \fi
  % set entry spacing
  \setlength{\itemsep}{#2\baselineskip}}}
 {\end{list}}
\usepackage{calc}
\newcommand{\CSLBlock}[1]{\hfill\break\parbox[t]{\linewidth}{\strut\ignorespaces#1\strut}}
\newcommand{\CSLLeftMargin}[1]{\parbox[t]{\csllabelwidth}{\strut#1\strut}}
\newcommand{\CSLRightInline}[1]{\parbox[t]{\linewidth - \csllabelwidth}{\strut#1\strut}}
\newcommand{\CSLIndent}[1]{\hspace{\cslhangindent}#1}
\usepackage{setspace} \usepackage{lineno} \usepackage{placeins} \usepackage[nottoc,notlof,notlot]{tocbibind} \renewcommand{\contentsname}{} \renewcommand{\listfigurename}{} \renewcommand{\listtablename}{} \usepackage{sectsty} \sectionfont{\centering}
\ifLuaTeX
  \usepackage{selnolig}  % disable illegal ligatures
\fi
\usepackage{bookmark}
\IfFileExists{xurl.sty}{\usepackage{xurl}}{} % add URL line breaks if available
\urlstyle{same}
\hypersetup{
  hidelinks,
  pdfcreator={LaTeX via pandoc}}

\author{}
\date{\vspace{-2.5em}}

\begin{document}

\doublespacing

\begin{center}
    
\textbf{\Large Applying effective environmental forecasts to Management Strategy Evaluation of yellowtail flounder}
    
\textsc{Sophie Wulfing$^{1*}$ Gavin Fay$^{1}$ Chengxue Li$^{3}$ Scott I. Large$^{2}$ Vincent Saba$^{2}$\\}
\vspace{3 mm}
\normalsize{\indent $^1$School for Marine Science and Technology, University of Massachusetts Dartmouth, University of New Hampshire, 02747, MA, USA\\ $^2$NOAA, National Marine Fisheries Service, Narragansett, RI 02882, USA\\ $^3$School of Marine and Atmospheric Sciences, Stony Brook University, Stony Brook, NY 11794, USA\\}
$\text{*}$ Corresponding authors: Sophie Wulfing (SophieWulfing@gmail.com)
\end{center}

\newpage

\linenumbers

\section{ABSTRACT}\label{abstract}

\section{INTRODUCTION}\label{introduction}

As the world's oceans are complex physical and chemical systems, climate change has impacted oceans in a variety of ways such increased ocean temperature and acidification, changes in ocean currents, (Fabry et al., 2008; Guinotte \& Fabry, 2008; Abraham et al., 2013; Sydeman et al., 2015; Wilson et al., 2016; Poloczanska et al., 2016; García Molinos et al., 2017; Friedland et al., 2022; Cheng et al., 2024). As a result, different marine species are projected to have complex and different responses to these environmental shifts depending on species' life histories, thermal tolerances, mobility, and other aspects of the species' biology. (Guinotte \& Fabry 2008; Chown et al 2010; Honisch et al 2012; Stewart et al 2014; Sydeman et al 2015; Polocanska et al 2016; Wilson et al 2016; Chaikin et al 2024). This creates compounding challenges when managing commercially-important fisheries. As some species will shift ranges northward or to deeper waters in search of colder temperatures, this will disrupt historical food chains as predators and prey may shift in different directions, or range shifts will introduce new interactions between predator and prey species that have historically not overlapped (Overholtz et al 2009; Albout et al 2014; Stewart et al 2014; Sydeman et al 2015; Wilson et al 2016; Garcia Molinos et al 2017; Chaikin et al 2024). Further, as fish populations may experience distribution shifts, annual growth patterns, or vary in abundance, this has potential to increase catch effort as historical fishing grounds and times may now not line up with fish stocks' new life history dynamics (Overholtz et al 2009; Poloczanska et al 2016; Tommasi 2017; Fulton et al 2021). New species may also shift their ranges into a commonly fished area, potentially opening up new commercial opportunities (Stewart et al., 2014).

\begin{itemize}
\item
  THIS IS WHAT HAPPENS IF THE KEY DOESN'T EXIST IN THE .BIB DOC: (\textbf{IfYouCiteAKeyThatDoesntExist?})
\item
  TO PUT A CITATION IN THE MIDDLE OF A SENTENCE, Sydeman et al. (2015) SHOWS YOU CAN JUST NOT INCLUDE THE BRACKETS
\item
  NOTICE HOW FORGETTING ONE @ SYMBOL OR BRACKET OR SEMI-COLON MESSES UP THE CITATION: {[}Fabry et al. (2008); Guinotte \& Fabry (2008); abrahamReviewGlobalOcean2013; Sydeman et al. (2015); Wilson et al. (2016); Poloczanska et al. (2016); García Molinos et al. (2017); Friedland et al. (2022); Cheng et al. (2024){]}
\end{itemize}

The US Northeast continental shelf has warmed faster than any other marine ecosystem in the United States (Saba et al 2016; Kleisner et al 2017; Saba et al 2023). The varying responses to climate change across fished species has presented challenges in accurately assessing stock health of various species in the region (Hare et al 2016; Kleisner et al 2017). Environmental covariates have rarely been incorporated into stock assessments in the Northeast (Mazure et al 2023). This has shown to reduce retrospective patterns and increase projection accuracy in New England populations such as Atlantic Cod, haddock, yellowtail flounder, and winter flounder (Buckely et al 2004; Klein et al 2017).

Management Strategy Evaluations (MSEs) are decision making frameworks that allow assessors to simulate and compare the sustainability and tradeoffs associated with different management decisions (Tommasi 2017; Mazure et al 2023; Saba et al 2023). Applying climate change covariates and MSE to new England groundfish has shown to improve the robustness of stock assessments (Mazur 2023). However, applications of MSEs can still be challenging, as climate change shifts population baselines and creates dynamics that are harder to predict in biological systems. This demonstrates a need for further research on how to incorporate these environmental conditions into the MSE framework (Hollowed et al 2011; Punt 2016; Holsman et al 2022; Saba et al 2023). Populations' responses to climate change can be complex, as fish exhibit different sensitivities to environmental conditions depending on their lifestage, and different environmental changes (i.e.~increased sea surface or bottom temperature, changes in salinity, and altered currents) will have different effects on populations (Hollowed et al 2011).

Further, decisions on how to project environmental conditions are important to consider when predicting stock dynamics. As assessments attempt to capture future population trends, forecasting environmental trends is also imperative, however this can be challenging as optimal decisions can change depending on scale, species, and which environmental covariates are used (Hollowed et al 2011; Tommasi 2017; Holsman et al 202).

\newpage

\textbf{References}

\phantomsection\label{refs}
\begin{CSLReferences}{1}{2}
\bibitem[\citeproctext]{ref-abrahamReviewGlobalOcean2013}
Abraham, J. P., Baringer, M., Bindoff, N. L., Boyer, T., Cheng, L. J., Church, J. A., Conroy, J. L., Domingues, C. M., Fasullo, J. T., Gilson, J., Goni, G., Good, S. A., Gorman, J. M., Gouretski, V., Ishii, M., Johnson, G. C., Kizu, S., Lyman, J. M., Macdonald, A. M., \ldots{} Willis, J. K. (2013). A review of global ocean temperature observations: {Implications} for ocean heat content estimates and climate change. \emph{Reviews of Geophysics}, \emph{51}(3), 450--483. \url{https://doi.org/10.1002/rog.20022}

\bibitem[\citeproctext]{ref-chengNewRecordOcean2024}
Cheng, L., Abraham, J., Trenberth, K. E., Boyer, T., Mann, M. E., Zhu, J., Wang, F., Yu, F., Locarnini, R., Fasullo, J., Zheng, F., Li, Y., Zhang, B., Wan, L., Chen, X., Wang, D., Feng, L., Song, X., Liu, Y., \ldots{} Lu, Y. (2024). New {Record Ocean Temperatures} and {Related Climate Indicators} in 2023. \emph{Advances in Atmospheric Sciences}, \emph{41}(6), 1068--1082. \url{https://doi.org/10.1007/s00376-024-3378-5}

\bibitem[\citeproctext]{ref-fabryImpactsOceanAcidification2008}
Fabry, V. J., Seibel, B. A., Feely, R. A., \& Orr, J. C. (2008). Impacts of ocean acidification on marine fauna and ecosystem processes. \emph{ICES Journal of Marine Science}, \emph{65}(3), 414--432. \url{https://doi.org/10.1093/icesjms/fsn048}

\bibitem[\citeproctext]{ref-friedlandMiddleAtlanticBight2022}
Friedland, K. D., Miles, T., Goode, A. G., Powell, E. N., \& Brady, D. C. (2022). The {Middle Atlantic Bight Cold Pool} is warming and shrinking: {Indices} from in situ autumn seafloor temperatures. \emph{Fisheries Oceanography}, \emph{31}(2), 217--223. \url{https://doi.org/10.1111/fog.12573}

\bibitem[\citeproctext]{ref-garciamolinosOceanCurrentsModify2017}
García Molinos, J., Burrows, M. T., \& Poloczanska, E. S. (2017). Ocean currents modify the coupling between climate change and biogeographical shifts. \emph{Scientific Reports}, \emph{7}(1), 1332. \url{https://doi.org/10.1038/s41598-017-01309-y}

\bibitem[\citeproctext]{ref-guinotteOceanAcidificationIts2008}
Guinotte, J. M., \& Fabry, V. J. (2008). Ocean {Acidification} and {Its Potential Effects} on {Marine Ecosystems}. \emph{Annals of the New York Academy of Sciences}, \emph{1134}(1), 320--342. \url{https://doi.org/10.1196/annals.1439.013}

\bibitem[\citeproctext]{ref-poloczanskaResponsesMarineOrganisms2016}
Poloczanska, E. S., Burrows, M. T., Brown, C. J., García Molinos, J., Halpern, B. S., Hoegh-Guldberg, O., Kappel, C. V., Moore, P. J., Richardson, A. J., Schoeman, D. S., \& Sydeman, W. J. (2016). Responses of {Marine Organisms} to {Climate Change} across {Oceans}. \emph{Frontiers in Marine Science}, \emph{3}. \url{https://doi.org/10.3389/fmars.2016.00062}

\bibitem[\citeproctext]{ref-stewartCombinedClimatePreymediated2014}
Stewart, J. S., Hazen, E. L., \& Bograd, S. J. (2014). Combined climate- and prey-mediated range expansion of {Humboldt} squid ({Dosidicus} gigas), a large marine predator in the {California Current System}. \emph{Global Change Biology}.

\bibitem[\citeproctext]{ref-sydemanClimateChangeMarine2015}
Sydeman, W. J., Poloczanska, E., Reed, T. E., \& Thompson, S. A. (2015). Climate change and marine vertebrates. \emph{Oceans and Climate}, \emph{350}.

\bibitem[\citeproctext]{ref-wilsonClimatedrivenChangesOcean2016}
Wilson, L. J., Fulton, C. J., Hogg, A. M., Joyce, K. E., Radford, B. T. M., \& Fraser, C. I. (2016). Climate-driven changes to ocean circulation and their inferred impacts on marine dispersal patterns. \emph{Global Ecology and Biogeography}, \emph{25}(8), 923--939. \url{https://doi.org/10.1111/geb.12456}

\end{CSLReferences}

\end{document}
